% ID: SEL-D0044-OB-E01
% ESTADO: validada. S045, cotejada con examen y criterios.
\textbf{a)} El recinto está a la derecha de $x=0$,
por encima de $y=(5x-42)/7$ y por debajo de
$y=3-2x/5$ y $y=6-x$. Los valores negativos de $y$
pueden pertenecer al recinto.
\begin{center}
\setlength{\unitlength}{5mm}
\begin{picture}(12,12)(-2,-7)
\qbezier(0,0)(4,0)(8,0)\put(8.2,0){$x$}
\qbezier(0,-6.5)(0,-1.4)(0,3.7)\put(0,3.8){$y$}
\qbezier(0,3)(2.5,2)(5,1)
\qbezier(5,1)(6,0)(7,-1)
\qbezier(7,-1)(3.5,-3.5)(0,-6)
\put(-0.2,3){\makebox(0,0)[r]{$(0,3)$}}
\put(-0.2,-6){\makebox(0,0)[r]{$(0,-6)$}}
\put(5.15,1.25){$(5,1)$}
\put(7.15,-1){$(7,-1)$}
\end{picture}
\end{center}
\textbf{b)} En $x=0$ se obtiene $-6\le y\le3$.
Las dos fronteras superiores se cortan en $(5,1)$ y la
inferior corta a $y=6-x$ en $(7,-1)$. El cruce de la inferior
con $y=3-2x/5$ queda por encima de $y=6-x$. Por tanto
los vértices son $(0,-6)$, $(0,3)$, $(5,1)$ y $(7,-1)$.

\textbf{c)} $F=3-2(x+y)$. Como $x+y\le6$, el mínimo es
$\boxed{-9}$. Se alcanza en todo el segmento entre $(5,1)$
y $(7,-1)$, donde $x+y=6$.
